\scanpage{44}\mspage{43}

$\mathcal F'$ est un \uline{sous-atelier} \uline{spacieux}, \unsure{si l'est|si c'est} une partie fermée \uline{spacieuse} de $\mathcal F$ \deleted{(\unsure{auquel cas} la condition Sousat 2 est} Ce qui suffit à entrainer que c'est un sous-atelier.

\uline{Exemple} Soit $\mathcal F$ un atelier, et considérons l'application canonique \unsure{:|.}
\[
\mathrm{Mulomb}\colon\mathcal F\hookrightarrow\mathrm{Fig}(M),
\]
on a vu que l'image est une partie fermée pour $\leq$, et que les relations $\leq$ et $\ll$ sur $\mathcal F$ sont induites par celles de $\mathrm{Fig}(M)$. Ainsi $\mathcal F$ est un sous-atelier de l'atelier des figures \unsure{dans} $M$. \deleted{En fait, on a vu aussi que c'est} \inserted{Mais ce n'est pas en général} un sous-atelier spacieux (mais ça l'est peut-être dans le cas ensembliste, à voir).

\par\noindent\hfill\rule{2cm}{.3pt}\hfill\mbox{}\par

\deleted{Je reviens à la \unsure{m} aux ``petites'' \unsure{ombres} de}

\uline{Remarque} Soit $\mathcal F$ un ordonné pour $\leq$, $\mathcal F'$ une \unsure{partie} fermée \inserted{non vide} de $\mathcal F$. Alors les parties fermées non vides de $\mathcal F'$ sont \unsure{exactmt|exactt} les parties fermées non vides de $\mathcal F$ qui sont contenues dans $\mathcal F'$. On a le même énoncé pour les parties fermées\unsure{.} spacieuses, pourvu que l'on suppose $\mathcal F'$ spacieux. De façon précise, si on a des
